% Generated by scripts/assemble.py; corrections belong in source/corrections.json.
\examhead{12}{北京科技大学 2025 年数学分析试卷}{17}
% p018-b12; source PDF page 18
\par\Needspace{5\baselineskip}
\blockanchor{p018-b12}
\majorquestion{一}解答如下问题. 每题 10 分, 共 40 分.\par
% p018-b13; source PDF page 18
\blockanchor{p018-b13}
\question{1}设数列 $a_{n} = \frac{\sum\limits_{k=0}^{n}\ln C_{n}^{k}}{n^{2}}$, 证明: $\lim_{n\to\infty}a_n = \frac{1}{2}$.\par
% p018-b14; source PDF page 18
\blockanchor{p018-b14}
\question{2}计算积分  $\int_{\frac{\pi}{6}}^{\frac{\pi}{3}}\frac{\cos^{2}x}{x(\pi-2x)}dx.$\par
% p018-b15; source PDF page 18
\par\Needspace{6\baselineskip}
\blockanchor{p018-b15}
\question{3}设两个数列 $\{a_{n}\}$ 与 $\{b_{n}\}$ 的定义如下:\par
\nopagebreak[4]
% p018-b16; source PDF page 18
\blockanchor{p018-b16}
\displayformula{a _ {1} = 3, b _ {1} = 2, a _ {n + 1} = a _ {n} + 2 b _ {n}, b _ {n + 1} = a _ {n} + b _ {n}, n = 1, 2, \dots .}
% p018-b17; source PDF page 18
\blockanchor{p018-b17}
而数列  $c_{n}=\frac{a_{n}}{b_{n}}, n=1,2,\cdots$ ，计算极限  $\lim_{n\to\infty}c_{n}$ .\par
% p018-b18; source PDF page 18
\blockanchor{p018-b18}
\question{4}讨论积分 $\int_{1}^{+\infty} e^{-\alpha x} \frac{\cos x}{x^p} \, \mathrm{d}x$ 在 $0 \leq \alpha < +\infty$ ($p > 0$ 是常数)的一致收敛性.\par
% p018-b19; source PDF page 18
\par\Needspace{5\baselineskip}
\blockanchor{p018-b19}
\majorquestion{二}(15 分) 若 $f(x), g(x)$ 在 $[0, 1]$ 上可积, 证明: $\lim_{n \to \infty} \frac{1}{n} \sum_{i=1}^{n} f\left(\frac{i}{n}\right) g\left(\frac{i-1}{n}\right) = \int_0^1 f(x) g(x) \, \mathrm{d}x$.\par
% p018-b20; source PDF page 18
\par\Needspace{5\baselineskip}
\blockanchor{p018-b20}
\majorquestion{三}(15 分) 证明不等式:\par
% p018-b21; source PDF page 18
\blockanchor{p018-b21}
(1) $\sqrt{t} - \frac{1}{\sqrt{t}} > \ln t (t > 1)$.\par
% p018-b22; source PDF page 18
\blockanchor{p018-b22}
(2) $\sum_{k=1}^{n} \frac{1}{\sqrt{k^2 + k}} > \ln(1 + n) (n \in \mathbb{N}^+)$.\par
% p018-b23; source PDF page 18
\par\Needspace{5\baselineskip}
\blockanchor{p018-b23}
\majorquestion{四}(10 分) 利用确界原理证明有限覆盖定理.\par
% p018-b24; source PDF page 18
\par\Needspace{5\baselineskip}
\blockanchor{p018-b24}
\majorquestion{五}(15 分) 证明:\par
% p019-b00; source PDF page 19
\blockanchor{p019-b00}
(1) 级数 $\sum_{n=1}^{\infty} \frac{2x^2}{1 + n^3x^4}$ 在 $[0, +\infty)$ 上一致收敛.\par
% p019-b01; source PDF page 19
\blockanchor{p019-b01}
(2) 级数 $\sum_{n=1}^{\infty} \frac{n^2}{1 + n^4 x^2}$ 在 $(0,1)$ 上收敛, 但非一致收敛.\par
% p019-b02; source PDF page 19
\par\Needspace{5\baselineskip}
\par\Needspace{6\baselineskip}
\blockanchor{p019-b02}
\majorquestion{六}(15 分) 若 $f'(x)$ 在 $[a, b]$ 上连续, 对任意正整数 $n$, 令 $x_i = a + \frac{b - a}{n} i, i = 1, 2, \cdots, n$, 记\par
\nopagebreak[4]
% p019-b03; source PDF page 19
\blockanchor{p019-b03}
\displayformula{A _ {n} = \frac {b - a}{n} \sum_ {i = 1} ^ {n} f (x _ {i}) - \int_ {a} ^ {b} f (x) \mathrm{d} x, n = 1, 2, \dots}
% p019-b04; source PDF page 19
\blockanchor{p019-b04}
证明： $\lim_{n\to \infty}nA_n = \frac{b - a}{2} [f(b) - f(a)].$\par
% p019-b05; source PDF page 19
\par\Needspace{5\baselineskip}
\blockanchor{p019-b05}
\majorquestion{七}(15 分) 设函数 $f(x, y)$ 在单位圆 $\{(x, y) \mid x^2 + y^2 = 1\}$ 上有连续的偏导数, 且 $f(0, 1) = f(1, 0)$, 证明: 在单位圆上至少有两点满足方程 $y \frac{\partial}{\partial x} f(x, y) = x \frac{\partial}{\partial y} f(x, y)$.\par
% p019-b06; source PDF page 19
\par\Needspace{5\baselineskip}
\par\Needspace{6\baselineskip}
\blockanchor{p019-b06}
\majorquestion{八}(15 分) 设 $\Omega$ 是不含原点的有界闭区域, 其体积为 $V$, 其边界 $S$ 为光滑的简单闭曲面, $\boldsymbol{n} = (\cos \alpha, \cos \beta, \cos \gamma)$ 为 $S$ 上点 $P(x, y, z)$ 处的单位外法向量. $\boldsymbol{r} = (x, y, z)$, $(\boldsymbol{n}, \boldsymbol{r})$ 为 $\boldsymbol{n}$ 与 $\boldsymbol{r}$ 的夹角, $f(t)$ 是 $[0, +\infty)$ 上的连续可微函数, 且满足 $tf'(t) + 2f(t) - t = 0$, 计算积分\par
\nopagebreak[4]
% p019-b07; source PDF page 19
\blockanchor{p019-b07}
\displayformula{\iint_ {S} f \left(\sqrt {x ^ {2} + y ^ {2} + z ^ {2}}\right) \cos (\boldsymbol {n}, \boldsymbol {r}) \mathrm{d} S.}
% p019-b08; source PDF page 19
\par\Needspace{5\baselineskip}
\par\Needspace{6\baselineskip}
\blockanchor{p019-b08}
\majorquestion{九}(10 分) 计算二重积分\par
\nopagebreak[4]
% p019-b09; source PDF page 19
\blockanchor{p019-b09}
\displayformula{\iint_ {D} (x \sin \alpha - y \cos \alpha) ^ {2} d x d y.}
% p019-b10; source PDF page 19
\blockanchor{p019-b10}
这里 $D = \left\{(x,y)\mid \frac{x^2}{a^2} +\frac{y^2}{b^2}\leq 1,a > 0,b > 0\right\} ,\alpha$ 为常数.\par
